\documentclass[12pt]{article}
\usepackage[T1]{fontenc}
\usepackage[utf8]{inputenc}
\usepackage{lmodern}
\usepackage{textcomp}
\usepackage{enumitem}
\usepackage{geometry}
\geometry{margin=1in}
\begin{document}
\begin{center}
Answer Key - Constitutional Law - Undergraduate Level Assessment 20 \
\small (Answers are ordered exactly as in the question paper.)
\end{center}

\section*{Multiple Choice}
\begin{enumerate}[label=\arabic*.]
\item \textbf{Answer:} C
\\ \textit{Options:} A) In hard cases judges generally decide cases on the basis of rights.; B) The rights of the parties feature in the determination of most cases before the courts.; C) Judges exercise strong discretion.; D) Judges seek the best 'fit' with constitutional and institutional history.
\\ \textit{Note:} Judges exercising strong discretion suggests that there may be multiple valid interpretations or outcomes in legal cases, contradicting Dworkin's assertion that there is only one right answer to every legal question.
\item \textbf{Answer:} A
\\ \textit{Options:} A) Doctrine of substituted security; B) Doctrine of due process; C) Doctrine of bias; D) Doctrine of reasonableness
\\ \textit{Note:} The Doctrine of Substituted Security is primarily associated with economic and property law rather than Natural Law thinking, which emphasizes inherent rights and moral principles derived from nature and human reason.
\item \textbf{Answer:} C
\\ \textit{Options:} A) Deterrence.; B) Rehabilitation.; C) Vengeance.; D) Desert.
\\ \textit{Note:} Durkheim viewed punishment primarily as a means of expressing societal outrage and reinforcing social cohesion, which aligns with the concept of vengeance as it reflects the collective response to transgressions against shared values.
\item \textbf{Answer:} C
\\ \textit{Options:} A) A society's law is a reflection of its culture.; B) Law is like language.; C) Law is the deliberate expression of a sovereign's will.; D) Law is an integral element of the social fabric.
\\ \textit{Note:} The correct answer is justified because Savigny's concept of Volksgeist emphasizes that law emerges organically from the unique cultural and historical context of a people rather than being a mere product of a sovereign's arbitrary will.
\item \textbf{Answer:} B
\\ \textit{Options:} A) Effectiveness of rehabilitation of multiplied by hazard to the community.; B) Extent of responsibility multiplied by actual harm done.; C) Risk of violence multiplied by degree of humility of offender.; D) Recidivism multiplied by defendant's history.
\\ \textit{Note:} The formula r x H reflects Robert Nozick's view that the appropriate punishment should be proportional to the extent of the offender's responsibility for the crime and the actual harm caused to the victim, emphasizing a fair and just approach to punishment in the context of constitutional law.
\end{enumerate}

\section*{True / False}
\begin{enumerate}[label=\arabic*.]
\setcounter{enumi}{5}
\item \textbf{Answer:} True
\\ \textit{Options:} A) True; B) False
\\ \textit{Note:} Kelsen's Grundnorm serves as a foundational principle that underpins the validity of a legal system, allowing us to comprehend the hierarchy and coherence of laws within that system.
\item \textbf{Answer:} False
\\ \textit{Options:} A) True; B) False
\\ \textit{Note:} Standard-setting in human rights diplomacy typically involves the development of non-binding guidelines and principles that aim to influence state behavior, rather than legally binding obligations.
\item \textbf{Answer:} True
\\ \textit{Options:} A) True; B) False
\\ \textit{Note:} Moral values provide a foundational framework for the development and interpretation of laws, influencing concepts of justice and establishing standards of right and wrong that are essential for a fair legal system.
\item \textbf{Answer:} False
\\ \textit{Options:} A) True; B) False
\\ \textit{Note:} The statement is false because the UN Human Rights Council is not a treaty-based mechanism; rather, it is an intergovernmental body established by the UN General Assembly that reviews the human rights records of member states through the Universal Periodic Review process.
\item \textbf{Answer:} True
\\ \textit{Options:} A) True; B) False
\\ \textit{Note:} The statement is true because under international law, the flag State has exclusive jurisdiction over crimes committed on its vessels, reflecting its sovereignty and responsibility to enforce laws on board.
\end{enumerate}

\section*{Long Form Response}
\begin{enumerate}[label=\arabic*.]
\setcounter{enumi}{10}
\item \textbf{Answer:} Treaties play a crucial role as sources of international law under Article 38 of the International Court of Justice (ICJ) Statute, particularly those that are in force and binding on the parties involved. This emphasis on treaties that are currently effective ensures that the obligations and rights established by the treaty are enforceable and recognized in international legal disputes. Treaties create legal commitments that parties must adhere to, thereby providing stability and predictability in international relations. Furthermore, the binding nature of these treaties underscores the principle of pacta sunt servanda, which obligates states to honor their agreements, enhancing the legitimacy and functionality of international law. Thus, only those treaties that are actively binding on the disputing parties can be effectively invoked in legal proceedings before the ICJ, reinforcing the treaty's significance in the resolution of international conflicts.
\\ \textit{Note:} correct because it underscores that treaties, as binding agreements recognized by parties, establish enforceable legal obligations that contribute to the stability and predictability of international law, as outlined in Article 38 of the ICJ Statute.
\item \textbf{Answer:} The exhaustion of local remedies in international human rights law serves a dual purpose: it seeks to alleviate the burden on international tribunals by limiting the number of petitions they receive, while simultaneously fostering the development and capacity of local judicial systems. By requiring individuals to first seek redress in their own courts, this principle encourages states to address human rights violations domestically, thereby strengthening local legal frameworks and promoting accountability. This process not only enhances the efficiency of international tribunals by prioritizing cases that have genuinely been unable to find resolution at the local level but also empowers local institutions to handle human rights issues effectively. Ultimately, the local remedies rule supports a more sustainable and equitable approach to human rights protection, ensuring that local courts are given the opportunity to fulfill their role.
\\ \textit{Note:} correct because it highlights that the exhaustion of local remedies not only reduces the caseload of international tribunals but also incentivizes states to improve their domestic judicial systems, thereby promoting accountability and enhancing the overall effectiveness of human rights protection.
\end{enumerate}

\vfill
\noindent\textit{End of Answer Key}
\end{document}